% 用 XeLaTeX 编译；公式、中文、图均在此处统一配置。
\usepackage[margin=25mm,headheight=14pt]{geometry}
\usepackage{fontspec}
\setmainfont{TeX Gyre Termes}
\setsansfont{TeX Gyre Heros}
\setmonofont{Latin Modern Mono}
\usepackage{amsmath,amssymb,mathtools,bm}
\usepackage{booktabs,tabularx,longtable,array}
\usepackage{graphicx}
\usepackage{xcolor}
\definecolor{BFblue}{HTML}{356A9A}
\definecolor{BFgreen}{HTML}{287C70}
\definecolor{BFviolet}{HTML}{7955A3}
\definecolor{BFamber}{HTML}{A97822}
\definecolor{BFink}{HTML}{243246}
\usepackage{tikz}
\usetikzlibrary{arrows.meta,positioning,calc,fit,backgrounds,shapes.geometric}
\tikzset{
  phys/.style={draw=BFblue,fill=BFblue!8,rounded corners=3pt,
  align=center,minimum height=9mm,inner sep=5pt,font=\small},
  nerve/.style={draw=BFviolet,fill=BFviolet!8,rounded corners=3pt,
  align=center,minimum height=9mm,inner sep=5pt,font=\small},
  therm/.style={draw=BFgreen,fill=BFgreen!8,rounded corners=3pt,
  align=center,minimum height=9mm,inner sep=5pt,font=\small},
  drive/.style={draw=BFamber,fill=BFamber!10,rounded corners=3pt,
  align=center,minimum height=9mm,inner sep=5pt,font=\small},
  arr/.style={-{Latex[length=2mm]},thick,draw=BFink},
  fb/.style={-{Latex[length=2mm]},semithick,dashed,draw=BFgreen}
}
\usepackage{enumitem}
\setlist{nosep,leftmargin=2em,itemsep=3pt}
\usepackage{caption}
\usepackage{listings}
\lstset{basicstyle=\ttfamily\small,frame=single,breaklines=true,
columns=fullflexible,keepspaces=true,showstringspaces=false,
numbers=left,numberstyle=\tiny,numbersep=6pt}
\usepackage[most]{tcolorbox}
\newtcolorbox{plainbox}{enhanced,breakable,colback=BFblue!4,
colframe=BFblue!35,boxrule=.5pt,arc=2pt}
\newtcolorbox{examplebox}{enhanced,breakable,colback=BFgreen!4,
colframe=BFgreen!40,boxrule=.5pt,arc=2pt}
\usepackage{hyperref}
\hypersetup{unicode=true,colorlinks=true,linkcolor=BFblue,
urlcolor=BFblue,pdftitle={BFZI 电磁热神经代理模型完整技术说明}}
\setlength{\parindent}{2em}
\setlength{\parskip}{3pt}
\linespread{1.18}
\newcommand{\zsea}{\mathbf Z_{\mathrm{sea}}}
\newcommand{\zcu}{\mathbf Z_{\mathrm{cu}}}
\newcommand{\ztot}{\mathbf Z_{\mathrm{EM}}}
\newcommand{\zext}{\mathbf Z_{\mathrm{ext}}}
\newcommand{\jj}{\mathrm j}
\newcommand{\R}{\mathbb R}
\newcommand{\C}{\mathbb C}
\newcommand{\dd}{\mathrm d}
\newcommand{\code}[1]{\texttt{#1}}
